\documentclass[conference]{IEEEtran}

% ------------------------------------------------
% Packages
% ------------------------------------------------
\usepackage{amsmath, amssymb}
\usepackage{graphicx}
\usepackage{cite}
\usepackage{url}
\usepackage{booktabs}
\usepackage{float}
\renewcommand{\refname}{Bibliografía}
% ------------------------------------------------
% Document
% ------------------------------------------------
\begin{document}

% ------------------------------------------------
% Title
% ------------------------------------------------
\title{Your Graph Theory Project Title}

% ------------------------------------------------
% Authors
% ------------------------------------------------
\author{
    \IEEEauthorblockN{
        Gabriel Salcedo,
        First Author [Last Name],
        Second Author [Last Name]
    }
    \IEEEauthorblockA{
        Escuela de Economía y Negocios\\
        Universidad del Rosario\\
        Bogotá, Colombia\\
        \{gabriel.salcedo, email2, email3\}@urosario.edu.co
    }
}

\maketitle


% ------------------------------------------------
% I. Introducción
% ------------------------------------------------
\section{Introducción}
Introduce el tema general de manera breve. Contextualiza la importancia de la teoría de grafos en esta área de estudio.

Al final de esta sección, debes explicar brevemente cómo está organizado el resto del documento: ``El resto de este documento está organizado de la siguiente manera. La Sección II detalla la descripción del problema...''

% ------------------------------------------------
% II. Descripción del Problema
% ------------------------------------------------
\section{Descripción del Problema}
Explica con detalle qué problema real o teórico estás abordando. ¿Qué está fallando o qué necesidad tecnológica/científica estás intentando resolver? Describe las limitaciones actuales del entorno.

% ------------------------------------------------
% III. Objetivos (General y Específicos)
% ------------------------------------------------
\section{Objetivo General y Objetivos Específicos}

\subsection{Objetivo General}
Escribe aquí el objetivo principal del proyecto. Debe ser un único enunciado claro que comience con un verbo en infinitivo (por ejemplo: \textit{Diseñar}, \textit{Optimizar}, \textit{Analizar}).

\subsection{Objetivos Específicos}
Pasos técnicos y secuenciales necesarios para alcanzar el objetivo general:
\begin{itemize}
    \item Primer paso técnico (ej. Recopilar los datos de la red).
    \item Segundo paso técnico (ej. Modelar la red mediante teoría de grafos).
    \item Tercer paso técnico (ej. Evaluar el rendimiento del algoritmo implementado).
\end{itemize}

% ------------------------------------------------
% IV. Marco Teórico
% ------------------------------------------------
\section{Marco Teórico}
Introduce los conceptos fundamentales de la teoría de grafos requeridos para entender tu proyecto.

\subsection{Conceptos Fundamentales}
Define matemáticamente los componentes del modelo. Por ejemplo, un grafo se puede representar formalmente como:
\begin{equation}
G = (V,E)
\end{equation}
donde $V$ representa el conjunto de vértices y $E$ representa el conjunto de aristas.

\subsection{Algoritmos o Métodos Utilizados}
Explica conceptualmente los teoremas o algoritmos específicos que aplicarás más adelante (ej. Algoritmo de Dijkstra, componentes conectadas, centralidad, etc.).

% ------------------------------------------------
% V. Modelamiento del Problema
% ------------------------------------------------
\section{Modelamiento del Problema}
Explica detalladamente cómo traduces el problema real de la Sección II a un modelo matemático de grafos basado en la Sección IV.

Abstrae los elementos del problema real de la siguiente forma:
\begin{itemize}
    \item \textbf{Vértices ($V$):} Representan [definición, ej. routers, ciudades, usuarios].
    \item \textbf{Aristas ($E$):} Representan [definición, ej. cables de fibra, carreteras, relaciones de amistad].
    \text{\textbf{Pesos de las aristas:}} Representan [definición opcional, ej. latencia, distancia en km].
\end{itemize}

% ------------------------------------------------
% VI. Solución Propuesta
% ------------------------------------------------
\section{Solución Propuesta}
Describe aquí la estrategia teórica, lógica o algorítmica elegida para resolver el problema ya modelado. Explica la arquitectura de tu propuesta, pseudocódigo o las reglas del sistema que vas a diseñar.

% ------------------------------------------------
% VII. Implementación de la Solución
% ------------------------------------------------
\section{Implementación de la Solución}
Detalla el desarrollo práctico real de tu propuesta. Explica las herramientas de software empleadas (ej. Python con NetworkX, MATLAB, C++), cómo procesaste los datos, librerías críticas o detalles de hardware si los hubiera.

% ------------------------------------------------
% VIII. Resultados Obtenidos
% ------------------------------------------------
\section{Resultados Obtenidos}
Presenta de forma cuantitativa y cualitativa los hallazgos de tus experimentos, simulaciones o ejecuciones de código.

% Ejemplo de Figura (Las leyendas van ABAJO en formato IEEE)
\begin{figure}[H]
    \centering
    \includegraphics[width=0.8\columnwidth]{graph.png} % Asegúrate de tener esta imagen en Overleaf o comenta esta línea
    \caption{Representación en grafo del modelo propuesto.}
    \label{fig:graph}
\end{figure}

Como se observa en la Fig.~\ref{fig:graph}, la topología de la red muestra las propiedades clave analizadas.

% Ejemplo de Tabla (Las leyendas van ARRIBA en formato IEEE con números romanos)
\begin{table}[H]
    \centering
    \caption{Propiedades del Grafo de la Red}
    \label{tab:properties}
    \begin{tabular}{lc}
        \toprule
        Propiedad & Valor \\
        \midrule
        Número de vértices & 10 \\
        Número de aristas & 15 \\
        Grado promedio & 3.0 \\
        Diámetro de la red & 4 \\
        \bottomrule
    \end{tabular}
\end{table}

En la Tabla~\ref{tab:properties} se resumen las métricas estructurales obtenidas tras la ejecución del algoritmo sobre la solución implementada.

% ------------------------------------------------
% IX. Conclusiones
% ------------------------------------------------
\section{Conclusiones}
Resume los logros principales del proyecto basándote en los resultados. Explica directamente cómo el uso de la teoría de grafos permitió solucionar el problema y menciona ideas puntuales para trabajos futuros o mejoras del sistema \cite{reference1}.

% ------------------------------------------------
% Bibliografía
% ------------------------------------------------
\bibliographystyle{IEEEtran}
\bibliography{references} % Recuerda tener un archivo llamado 'references.bib' en tu proyecto

\end{document}
